\begin{remark}[7.9.6]
\label{IV.2.s7.9.6-fr}
Es posible que el recíproco de (7.9.5) sea verdadero, es decir, que
la hipótesis de que $X$ sea reducido y de que los anillos de los abiertos afines de $X$
satisfagan las condiciones (ii) y (iii) de (7.8.2) implique la posibilidad de
resolver $X$ (y, por tanto, también la de resolver todo preesquema
reducido localmente de tipo finito sobre $X$, en virtud de (7.4.4) y (6.12.4)).
Esto es cierto, en todo caso, si solo se consideran los preesquemas \emph{noetherianos} reducidos
cuyos cuerpos residuales sean de \emph{característica} $0$, como muestran
los resultados recientes de Hironaka [35] (este último enuncia sus resultados bajo
hipótesis demasiado restrictivas, pues sus razonamientos son en realidad válidos
bajo las solas hipótesis (ii) y (iii) de (7.8.2) para los anillos de los
abiertos afines de $X$, junto con el hecho de que los cuerpos residuales sean de
característica $0$). Por el contrario, en el caso general, hasta ahora
sólo se sabe resolver los esquemas algebraicos de dimensión $2$ sobre un cuerpo perfecto
[36]. Dada la importancia que está adquiriendo la resolución de singularidades
en el estudio topológico de los esquemas (en particular, por lo que se refiere
a sus propiedades homológicas y homotópicas), esto confiere un interés
especial a la categoría de los anillos excelentes o de los preesquemas
excelentes.

También cabe observar a este respecto que la parte menos delicada de los
razonamientos de Hironaka \emph{(loc. cit.)} muestra que, con independencia de toda
hipótesis sobre las características de los cuerpos residuales, y utilizando
únicamente las propiedades (ii) y (iii) de (7.8.2) para los anillos locales
de un preesquema $X$, el problema de la resolución de singularidades de $X$ se
reduce al caso en que $X=\operatorname{Spec}(A)$, siendo $A$ un anillo local
noetheriano íntegro y \emph{completo}. Por consiguiente, si resultados
posteriores llegasen a refutar la conjetura formulada más arriba y,
por tanto, obligasen a imponer condiciones más restrictivas a los
anillos locales de $X$, difícilmente podría hacerse de otro modo que mediante
condiciones restrictivas impuestas a anillos locales completos (probablemente
condiciones relativas a sus cuerpos residuales, tal vez la condición
$[k:k^p]<+\infty$ ($p$ exponente característico de $k$)).
\end{remark}

\par\smallskip
{\small\noindent\textbf{Nota histórica de la edición.}
Las expresiones temporales y las conjeturas de esta observación reflejan el estado de la cuestión en la publicación original de 1965; no describen el estado actual de la investigación.}\par
